\section{Discounted Cash Flow Analysis}

We forecast unlevered free cash flow explicitly for FY2026E--FY2030E and capitalize the residual value beyond FY2030E using a Gordon growth terminal value. Our approach follows standard practice for CAPM-based cost of capital and DCF construction \citep{damodaran2012,koller2020}.

\subsection{Weighted Average Cost of Capital}

We estimate a WACC of 9.8\%, built from a CAPM cost of equity and an after-tax cost of debt weighted at Vertex's targeted capital structure (Table~\ref{tab:wacc}). The equity risk premium and beta inputs follow standard market-practice sourcing \citep{sharpe1964,fama1993,duffphelps2026}.

\begin{table}[H]
\centering
\small
\caption{WACC build}
\label{tab:wacc}
\begin{tabular}{@{}lr@{}}
\toprule
\textbf{Component} & \textbf{Value} \\
\midrule
Risk-free rate (10-year sovereign benchmark) & 4.2\% \\
Equity risk premium & 5.5\% \\
Levered beta & 1.15 \\
Cost of equity (CAPM) & 10.5\% \\
\addlinespace
Pre-tax cost of debt & 5.8\% \\
Marginal tax rate & 24.0\% \\
After-tax cost of debt & 4.4\% \\
\addlinespace
Target equity / total capitalization & 88\% \\
Target debt / total capitalization & 12\% \\
\midrule
\textbf{WACC} & \textbf{9.8\%} \\
\bottomrule
\end{tabular}
\end{table}

\subsection{Free Cash Flow Projection}

Revenue growth decelerates from 22.0\% in FY2026E to 9.0\% in FY2030E as the Company laps its largest cohort of enterprise wins and the addressable base matures; EBITDA margin expands from 28.5\% to 34.0\% over the same period as Managed Cloud Services scales and professional-services intensity declines. Capital expenditure eases from 5.5\% to 4.5\% of revenue as platform infrastructure investment plateaus, and incremental net working capital investment is held at 1.5\% of revenue. Table~\ref{tab:dcf-fcf} presents the full build.

\begin{table}[H]
\centering
\small
\caption{Unlevered free cash flow projection (\$ in millions)}
\label{tab:dcf-fcf}
\begin{tabular}{@{}lrrrrr@{}}
\toprule
& \textbf{FY2026E} & \textbf{FY2027E} & \textbf{FY2028E} & \textbf{FY2029E} & \textbf{FY2030E} \\
\midrule
Revenue & 1,513 & 1,800 & 2,088 & 2,339 & 2,550 \\
\quad Revenue growth (\%) & 22.0\% & 19.0\% & 16.0\% & 12.0\% & 9.0\% \\
EBITDA & 431 & 540 & 658 & 772 & 867 \\
\quad EBITDA margin (\%) & 28.5\% & 30.0\% & 31.5\% & 33.0\% & 34.0\% \\
Less: depreciation \& amortization & (68) & (76) & (84) & (94) & (102) \\
\midrule
EBIT & 363 & 464 & 574 & 678 & 765 \\
Less: cash taxes @ 24.0\% & (87) & (111) & (138) & (163) & (184) \\
\midrule
NOPAT & 276 & 353 & 436 & 515 & 581 \\
Plus: depreciation \& amortization & 68 & 76 & 84 & 94 & 102 \\
Less: capital expenditures & (83) & (94) & (102) & (110) & (115) \\
Less: increase in net working capital & (23) & (27) & (31) & (35) & (38) \\
\midrule
\textbf{Unlevered free cash flow} & \textbf{238} & \textbf{308} & \textbf{387} & \textbf{464} & \textbf{530} \\
Discount factor @ 9.8\% (mid-year) & 0.954 & 0.869 & 0.792 & 0.721 & 0.657 \\
\textbf{Present value of free cash flow} & \textbf{227} & \textbf{268} & \textbf{306} & \textbf{335} & \textbf{348} \\
\bottomrule
\end{tabular}
\end{table}

\begin{figure}[H]
\centering
\begin{tikzpicture}
\begin{axis}[
  ybar,
  bar width=15pt,
  width=0.9\linewidth,
  height=4.6cm,
  enlarge x limits=0.15,
  ymin=0,
  ymax=950,
  ylabel={\$ in millions},
  ylabel style={font=\small,color=slateink},
  symbolic x coords={FY2026E,FY2027E,FY2028E,FY2029E,FY2030E},
  xtick=data,
  tick label style={font=\small,color=slateink},
  axis line style={color=slateink},
  ymajorgrids,
  grid style={gray!25},
  legend style={at={(0.5,-0.24)},anchor=north,legend columns=2,draw=none,font=\small},
]
\addplot[fill=meridiannavy,draw=none] coordinates {(FY2026E,431)(FY2027E,540)(FY2028E,658)(FY2029E,772)(FY2030E,867)};
\addplot[fill=vertexteal,draw=none] coordinates {(FY2026E,238)(FY2027E,308)(FY2028E,387)(FY2029E,464)(FY2030E,530)};
\legend{EBITDA, Unlevered free cash flow}
\end{axis}
\end{tikzpicture}
\caption{Projected EBITDA and unlevered free cash flow, FY2026E--FY2030E (\$ in millions). The widening gap reflects cash taxes, capital expenditure, and working-capital investment absorbing a roughly constant share of EBITDA even as the margin itself expands.}
\label{fig:fcf-ebitda}
\end{figure}

\subsection{Terminal Value and Enterprise Value}

We apply a 3.0\% terminal growth rate, modestly above long-run nominal GDP growth expectations and consistent with Vertex's exposure to a still-expanding data-infrastructure category. Table~\ref{tab:dcf-summary} bridges from the present value of explicit-period cash flow to an implied value per share.

\begin{table}[H]
\centering
\small
\caption{DCF enterprise and equity value bridge (\$ in millions, except per-share data)}
\label{tab:dcf-summary}
\begin{tabular}{@{}lr@{}}
\toprule
\textbf{Item} & \textbf{Value} \\
\midrule
Sum of PV of FY2026E--FY2030E unlevered FCF & 1,484 \\
FY2031E terminal-year FCF (FY2030E $\times$ 1.03) & 546 \\
Terminal value (FY2031E FCF / (WACC $-$ $g$)) & 8,029 \\
Discount factor @ 9.8\%, year 5 (end-of-year) & 0.626 \\
Present value of terminal value & 5,023 \\
\midrule
\textbf{Implied enterprise value} & \textbf{6,507} \\
Less: net debt & (240) \\
\midrule
\textbf{Implied equity value} & \textbf{6,267} \\
Diluted shares outstanding (mm) & 142.5 \\
\midrule
\textbf{Implied value per share} & \textbf{\$44.00} \\
\bottomrule
\end{tabular}
\end{table}

As a cross-check, the implied terminal value corresponds to an exit multiple of 9.3x FY2030E EBITDA -- meaningfully below the 14.0x--16.0x range at which the comparable-company set in Section~4 currently trades, which we view as an appropriately conservative discount given the deceleration to 9.0\% growth embedded in the terminal year.

\subsection{Sensitivity Analysis}

Table~\ref{tab:sensitivity} shows implied value per share across a plausible range of WACC and terminal growth assumptions. Our base case (WACC 9.8\%, $g$ 3.0\%) sits near the center of the grid.

\begin{table}[H]
\centering
\small
\caption{Sensitivity of implied value per share to WACC and terminal growth rate}
\label{tab:sensitivity}
\begin{tabular}{@{}lccc@{}}
\toprule
& \multicolumn{3}{c}{\textbf{Terminal growth rate ($g$)}} \\
\cmidrule(l){2-4}
\textbf{WACC} & \textbf{2.5\%} & \textbf{3.0\%} & \textbf{3.5\%} \\
\midrule
9.3\% & \$45.80 & \$47.65 & \$49.75 \\
9.8\% (base case) & \$42.20 & \cellcolor{signalgold!20}\textbf{\$44.00} & \$46.05 \\
10.3\% & \$38.95 & \$40.55 & \$42.35 \\
\bottomrule
\end{tabular}
\end{table}
